\section{Keys to Spiritual Development}

\begin{quotex}
If we see further than they, it is not in virtue of our stronger sight, but because we are lifted up by them and carried to great height. We are dwarfs carried on the shoulders of giants. \flright{\textsc{Bernard of Chartres}}

\end{quotex}
The  Cathedral of Our Lady of Chartres of has long been a destination for pilgrims and for esoteric study. \textbf{John of Salisbury}\footnote{\url{https://plato.stanford.edu/entries/john-salisbury/}} was an English priest, philosopher, diplomat, and eventually Bishop of Chartres from 1176 to 1180. In his works he describes seven keys of spiritual development and five levels of consciousness indicating the journey that the soul makes each night. These are worth study, if you want to kick up your game for 2019.

\paragraph{Seven Keys}
John of Salisbury expanded on the keys originally proposed by Bernard of Chartres.

\textbf{Humility of Mind}. Humility acquiesces in what the better informed have written.

\textbf{Zeal in the quest}. Because with burning zeal of love [the Apostles] were seeking what lay hidden, wisdom which was being sought was present and opened their eyes so that they understood the Scriptures, and they reproached their own slowness when they saw the Lord. For zealous investigation is profitable only when the desire for knowledge has as its end Christ.

\textbf{A serene life}. If zeal is, as Cicero declares, an ardent turning of the mind to the accomplishment of something which gives great pleasure, naturally a troubled mind torn in different directions is by no means devotedly occupied with the sole task of virtue. A peaceful life is a necessity even the moralist [Cato] teaches, for he thinks that the training of the scholar is unavailing if a crowding noisy throng is hammering at his door. This serenity is an inner state.

\textbf{Silence}. Zealous application is of the greatest help when virtue is reinforced as the result of a silent exercise of his own best judgment in respect to all that a man reads or hears; for there reason ponders and weighs in the balance the benefit derived from all things.

\textbf{Poverty}. A tranquil existence however is impossible if it lack the necessities of life; and again, if the mind luxuriates in pleasure, affluence extinguishes the light of reason. Consequently poverty, the breeder of men, the guard of true humility, and the ally of virtue is added to the preceding keys of philosophy that it, with its spirit of moderation, may check luxury and by its incentive force man to remember who he is and oblige him to press on without interruption in his task of virtue.

\textbf{A foreign land}. To think beyond one's circumstances because ``everywhere abides with wisdom". Philosophy exacts sojourn in a foreign land — nay, makes at times her own a foreign land — nor does she ever feel the burden of exile. This is because she drives away domestic worries, which are of the flesh, with the result that the man entirely devoted, as it were, to the spirit, regards as foreign all that impedes the progress of wisdom. Everywhere he is at home and everywhere in his own country, because everywhere abides with him and everywhere abides with wisdom.

\textbf{Love for one's teacher}. Instructors are to be loved and respected as parents are; for as the latter are creators of the bodies, so the former are the creators of the souls of their listeners. Such attachments are of great assistance to study, for pupils are glad to listen to those whom they love; they believe what they are told by them, desire to be like them; under the impulse of loyalty and affection they are eager and glad to form those throngs of pupils; do not become angry when rebuked; are not confused when praised, and will themselves well deserve to be held very dear because of their devotion to study. For as the function of teachers is to teach, so that of auditors to show themselves teachable; otherwise neither without the help of the other avails. Just as man owes his origin to each of two progenitors and as you labor in vain if you scatter seed on ground that has not been broken and softened by the furrow, so eloquence cannot mature unless there be a spirit of harmony between the teacher and the taught.

\paragraph{Journey of the Soul}
John of Salisbury develops this idea from Macrobius, who described five types of dreams.

\begin{enumerate}
\item \textbf{Insomnium}. Troubled dream cause by physical or mental disturbance. Troubled dreams are in general the result of insobriety or drunkenness, different emotions, turmoil of feeling, or vestiges of thoughts. The frenzied souls of lovers are always invaded by troubled dreams. 
\item \textbf{Phantasma}. An hallucination, a dream on the borderland of a state of being asleep and awake, when one sees strange vague forms flitting around him, to which type belongs the nightmare. In a nightmare, a person, as the result of various kinds of oppression, as though in semi-wakefulness or restless sleep imagines himself to be awake when in reality he is asleep, and feels himself crushed down by someone. 
\item \textbf{Somnium}. An ordinary dream contains images of events wrapped in a cloak of disguise, and it is with this disguise that the art of dream interpretation deals. 
\item \textbf{Oraculum}. A prophetic dream, in which some person of eminence appears — priest or even god — who delivers oracular pronouncements. When a communication is made in sleep by the agency of a second person and this individual is of honourable position, worthy of reverence, we have the oracular dream. An oracle is the pronouncement of divine will by the mouth of man or anything that assumes the human form: a human being, angel, god, or what you will. Moreover, an individual is honourable and worthy of reverence either as the result of nature as in the case of a parent, of position as in the case of a master, of character as in a man of piety, of chance as in a magistrate, of religion as in a god, angel, or man consecrated to holy office. 
\item \textbf{Visio}. A vision which produces in exact detail an event that is to occur in the future. This state is to be wide awake in the spiritual world and able to see one's surroundings quite clearly. 
\end{enumerate}
\textbf{Rene Querido} asserts that the stages of somnium, oraculum, and visio are actually stages of initiation.



References:

\textbf{Rene Querido}, \emph{The Golden Age of Chartres}

\textbf{John of Salisbury}, \emph{Policraticus}

\textbf{John of Salisbury}, \emph{Metalogicon}



\flrightit{Posted on 2019-01-01 by Cologero }
